\section*{Bab \ref{ch:b2:global-change} --- Perubahan Global dan Biosfer}

\begin{solution}{exo:b2:global-change:1}
$\Delta T = \lambda C_{\text{kum}}$ dengan $\lambda = \qty{1.65}{\celsius}$
tiap \qty{1000}{GtC}: \qty{1.2}{\celsius}, \qty{1.65}{\celsius},
\qty{3.3}{\celsius}.
\end{solution}

\begin{solution}{exo:b2:global-change:2}
Derajat-hari: jumlah kelebihan suhu rerata harian di atas sebuah dasar
sepanjang harinya, yaitu \omterm{thm:b2:global-change:phenology}{waktu termal} yang mengatur perkembangan
tumbuhan dan hewan berdarah dingin. Ketakserasian: ketika dua spesies
yang berinteraksi mengatur waktu musimnya dengan isyarat yang berbeda
dan pemanasannya memindahkan yang satu dan bukan yang lain --- burung
sikatannya, yang berisyarat panjang hari di Afrika, tiba pada
tanggalnya yang lama lalu mendapati puncak ulatnya, yang berisyarat
suhu, sudah lewat dua pekan.
\end{solution}

\begin{solution}{exo:b2:global-change:3}
Kira-kira \qty{150}{m} ke atas bukit dan \qty{150}{km} ke arah kutub
tiap derajatnya. Sebuah spesies di puncaknya tidak memiliki \omterm{def:b2:living-soil:soil}{tanah} yang
lebih tinggi: isotermnya naik di atas gunungnya dan habitatnya lenyap,
sementara luas tiap pitanya menyusut bersama ketinggiannya dalam
perjalanan ke atas.
\end{solution}

\begin{solution}{exo:b2:global-change:4}
\ensuremath{\mathrm{CO_2}} yang terlarut membentuk asam karbonat, yang
melepaskan ion hidrogen, dan pH-nya turun sebanding dengan logaritma
tekanan \ensuremath{\mathrm{CO_2}}-nya. Ion hidrogen tambahannya mengubah
karbonat menjadi bikarbonat, sehingga kepekatan karbonatnya dan \omterm{thm:b2:global-change:acid}{keadaan
kejenuhan} kalsium karbonatnya turun. Jasad yang membangun aragonit ---
karang, pteropoda, banyak larva --- dan kalsit --- kokolitofor,
foraminifera, moluska --- harus membelanjakan lebih banyak untuk
membangun cangkangnya dan, di bawah kejenuhannya, kehilangannya.
\end{solution}

\begin{solution}{exo:b2:global-change:5}
\qty{1.5}{\celsius}: $1.5/1.65\times 1000 = \qty{909}{GtC}$ seluruhnya,
dengan \qty{209}{GtC} tersisa, yaitu 21 tahun. \qty{2}{\celsius}:
\qty{1212}{GtC}, dengan \qty{512}{GtC} tersisa, yaitu 51 tahun.
\end{solution}

\begin{solution}{exo:b2:global-change:6}
$\sqrt{2\times 150/0.12} = 50$ hari; majunya $1.5/0.12 = 12.5$
hari.
\end{solution}

\begin{solution}{exo:b2:global-change:7}
\qty{3}{\celsius}: isotermnya naik $3000/6.5 = \qty{460}{m}$; pitanya
berpindah ke \qtyrange{1260}{1860}{m}, tetapi gunungnya berakhir pada
\qty{1800}{m}: \qty{60}{m} teratas pitanya hilang dan sisanya duduk di
tempat luasnya $2^{-460/300} = 0.35$ dari yang dahulu --- yaitu
kira-kira dua pertiga luasnya lenyap. \qty{5}{\celsius}: \qty{770}{m};
pitanya akan berada pada \qtyrange{1570}{2170}{m}; hanya
\qtyrange{1570}{1800}{m} yang ada, yaitu sepertiga lebar pitanya pada
$2^{-770/300} = 0.17$ luasnya --- di bawah sepersepuluh luas asalnya:
spesiesnya pada dasarnya lenyap.
\end{solution}

\begin{solution}{exo:b2:global-change:8}
560: pH $7.93$, ion hidrogennya $\times 1.9$; 800: $7.79$, $\times
2.6$; 1000: $7.70$, $\times 3.1$.
\end{solution}

\begin{solution}{exo:b2:global-change:9}
$1 - (1 - h)^{0.25}$: \qty{26}{\%}, \qty{53}{\%}, \qty{82}{\%}.
Kehilangan seketikanya lebih kecil karena sisanya masih menahan
populasi sebagian besar spesiesnya; populasinya terlalu kecil untuk
bertahan lalu mati sepanjang berpuluh tahun --- yaitu \omterm{thm:b2:global-change:sar}{utang kepunahan}.
\end{solution}

\begin{solution}{exo:b2:global-change:10}
$c = 2\sqrt{0.7\times 20} = \qty{7.5}{km/yr}$; \qty{1000}{km} dalam
kira-kira 130 tahun. Menyetengahkan $r$ atau $D$ membagi $c$ dengan
$\sqrt 2$, menjadi \qty{5.3}{km/yr}; tak satu pun menyetengahkannya ---
untuk menyetengahkan kecepatannya orang harus memperempat hasil kali $rD$.
\end{solution}

\begin{solution}{exo:b2:global-change:11}
Latarnya: $8\times 10^{6}\times 10^{-6} = 8$ tiap tahun, 800 tiap abad;
pada seribu kali lipatnya, 8000 tiap tahun, \num{800000} tiap abad.
Sepuluh ribu yang terdokumentasi dalam satu abad adalah kira-kira
sepuluh kali latarnya, bukan seribu kali --- tetapi dokumentasinya
hanya meliputi beberapa persen spesiesnya yang terperikan dan terpantau
(burung, mamalia, beberapa tumbuhan), yang di antaranya lajunya memang
seratus kali latarnya atau lebih; bagi mayoritasnya yang belum
terperikan, sebagian besar kepunahannya tidak pernah terlihat, dan
taksiran spesies--luasnya adalah satu-satunya pegangan. Kebenarannya
terletak antara yang tercacah dan yang tertaksir, dan itulah sebabnya
jangkauan yang dikutip begitu lebar.
\end{solution}

\begin{solution}{exo:b2:global-change:12}
Suhunya sebanding dengan \omterm{prop:b2:global-change:tcre}{pancaran kumulatifnya}; ketika pancaran
bersihnya nol, jumlah kumulatifnya berhenti bertambah dan begitu pula
pemanasannya. Di dalam \omterm{thm:b2:biogeochemical-cycles:box}{model kotaknya}, dengan pancarannya nol,
kelebihan atmosfernya mulai turun ketika lubuknya terus bekerja, yang
akan mendinginkan --- tetapi lautannya, yang masih menghangat menuju
kesetimbangan dengan \ensuremath{\mathrm{CO_2}} yang dinaikkan, terus
menghangatkan permukaannya, dan kedua efeknya nyaris saling meniadakan:
suhunya tetap berada di tempatnya selama berabad-abad. Yang hilang dari
hujahnya: gas rumah kaca lainnya (umur metana yang pendek berarti
pemanasannya turun cepat setelah pancarannya berhenti), aerosol yang
kini menutupi sebagian pemanasannya lalu lenyap dalam beberapa pekan
setelah pancaran yang membuatnya, dan umpan baliknya --- karbon lapisan
beku abadi, matinya hutan --- yang dapat menambahkan pancarannya
sendiri.
\end{solution}

\begin{solution}{pb:b2:global-change:1}
\textbf{1.} A: $700 + 80\times 10 = \qty{1500}{GtC}$. B: $700 +
\tfrac12\times 50\times 10 = \qty{950}{GtC}$.
\textbf{2.} A: $1.65\times 1.5 = \qty{2.5}{\celsius}$. B:
$1.65\times 0.95 = \qty{1.6}{\celsius}$.
\textbf{3.} Tahun 2050. A: \qty{1000}{GtC}, \qty{1.65}{\celsius}. B:
pancarannya $10\int_0^{30}(1 - t/50)\,\mathrm{d}t = 10(30 - 9) =
\qty{210}{GtC}$, seluruhnya 910, \qty{1.5}{\celsius}.
\textbf{4.} $\lambda\times 100 = \qty{0.165}{\celsius}$ tiap
dasawarsa, sedikit di bawah $0.2$ yang teramati (yang mencakup gas
lainnya dan hilangnya penutupan aerosolnya).
\textbf{5.} Ya: \omterm{prop:b2:global-change:tcre}{pancaran kumulatifnya} berhenti bertambah pada tahun
2070, dan begitu pula pemanasannya, yang tetap pada \qty{1.6}{\celsius}
--- kesebandingannya berarti suhunya ditetapkan jumlahnya, bukan
lajunya.
\textbf{6.} $2/1.65\times 1000 = \qty{1212}{GtC}$; jalur A mencapainya
setelah $512/10 = 51$ tahun, yaitu pada tahun 2071.
\textbf{7.} $\sqrt{2\times 200/0.1} = 63$ hari.
\textbf{8.} Pemanasan terhadap tahun 2020: A $2.5 - 1.2 =
\qty{1.3}{\celsius}$, B \qty{0.4}{\celsius}. Majunya $\Delta T/b$:
A 13 hari, B 4 hari.
\textbf{9.} Celahnya: A $20 + 13 = 33$ hari; B 24 hari.
\textbf{10.} Ulatnya tersedia dari hari 15 sampai hari 25 setelah
pendaunannya. B: burungnya tiba pada hari 24, yang masih di dalam
jendelanya, nyaris saja. A: hari 33, yaitu delapan hari setelah ulat
terakhirnya.
\textbf{11.} Burungnya maju $3\times 1.3 = 4$ hari: celahnya $33 - 4 =
29$ hari --- masih di luar jendelanya.
\textbf{12.} Burungnya dapat hidup pada suhu yang baru; yang tidak
dapat ia kerjakan adalah memberi makan anaknya ketika makanannya telah
lenyap, karena isyaratnya dan isyarat mangsanya menanggapi pemanasannya
secara berbeda.
\textbf{13.} A: $1.3/6.5\times 1000 = \qty{200}{m}$ ke atas dan
$1.3/0.65\times 100 = \qty{200}{km}$ ke arah kutub. B: \qty{57}{m} dan
\qty{57}{km}.
\textbf{14.} Dalam 80 tahun pohonnya berpindah \qty{8}{km}: ia tidak
dapat mengikuti \qty{200}{km} (A) dan juga kurang dari \qty{57}{km} (B)
--- pohon tertinggal, dan sebarannya akan tidak setimbang dengan
iklimnya selama berabad-abad pada jalur mana pun.
\textbf{15.} A: pitanya akan berada pada \qtyrange{1700}{2000}{m}: ia
tepat mencapai puncaknya, tanpa tempat lagi untuk pergi. Gunungnya
menyempit bersama ketinggiannya, sehingga irisan yang sama hanya
meliputi $2^{-200/300} = 0.63$ kali luasnya yang dahulu --- dan
pemanasan selanjutnya tidak memiliki \omterm{def:b2:living-soil:soil}{tanah} yang tersisa untuk
memindahkannya.
\textbf{16.} B: \qty{57}{m} ke atas, dengan faktor luas $2^{-57/300} = 0.88$: yaitu sebuah kehilangan \qty{12}{\%}.
\textbf{17.} Delapan dasawarsa pada \qty{30}{km}: \qty{240}{km}, lebih
banyak daripada \qty{200}{km} yang dituntut pada A: ikannya
mengimbangi (lalu berpindah keluar dari perikanan lamanya).
\textbf{18.} Pada tepi bawahnya spesiesnya berjumpa dengan pesaing,
patogen, dan pemangsa yang bergerak naik dari bawah yang belum pernah
ia hadapi, dan faalnya sendiri, yang teradaptasi dengan suhu yang lama,
dilampaui faal mereka: mundurnya bersifat persaingan sebanyak
bersifat faali.
\textbf{19.} A: $\mathrm{pH} = 8.2 - 0.9\log_{10}(600/280) = 7.90$;
ion hidrogennya $10^{0.30} = 2.0$ kali tahun 1750. B: $8.04$; $1.44$
kali.
\textbf{20.} $\Omega = 4/(\text{nisbahnya})$: A $2.0$, B $2.8$. Pada 2
karangnya masih mengapur tetapi dengan lambat, dengan rangkanya lebih
lemah dan terumbunya terkikis lebih cepat daripada ia tumbuh; pada
$2.8$ ia membangun, walaupun panas jalur B masih memutihkannya.
\textbf{21.} A: $400\times 0.4^{0.25} = 318$, dengan 82 hilang. B:
$400\times 0.8^{0.25} = 378$, dengan 22 hilang.
\textbf{22.} A: $318\times(1 - 0.05\times 1.3) = 298$. B:
$378\times(1 - 0.05\times 0.4) = 371$.
\textbf{23.} A: $1 - 298/400 = \qty{26}{\%}$. B: \qty{7}{\%}.
\textbf{24.} Pada keduanya, pembukaan lahannya: 82 berbanding 20
spesies pada A, 22 berbanding 7 pada B --- suku pemanasan modelnya
sederhana; dalam kenyataan keduanya berinteraksi, karena sebuah hutan
yang terserpih tidak dapat menggeser sebarannya.
\textbf{25.} Jalur A: \qty{2.5}{\celsius}, pH $7.90$, dengan seperempat
spesies pohonnya hilang. Jalur B: \qty{1.6}{\celsius}, pH $8.04$,
dengan \qty{7}{\%} hilang.
\end{solution}
